\documentclass[11pt,a4paper]{article}
\usepackage{eumine_hackathon}

%% -- Document metadata ------------------------------------------------------
\hypersetup{
  pdftitle={Database Access Guide --- EuMINe DataBridge Hackathon 2026},
  pdfauthor={EuMINe Hackathon Organizing Committee},
  pdfsubject={How to access Materials Project, JARVIS-DFT, AFLOW, NOMAD, OQMD},
}

\title{%
  {\large \textcolor{eumine_blue}{\textsc{EuMINe} DataBridge Hackathon 2026}}\\[6pt]
  \textbf{Database Access Guide}\\[4pt]
  {\normalsize Materials Project \textbullet\ JARVIS-DFT \textbullet\ AFLOW
   \textbullet\ NOMAD \textbullet\ OQMD}
}
\author{EuMINe Hackathon Organizing Committee}
\date{May 2026 --- v1.0}

\begin{document}
\maketitle
\thispagestyle{fancy}
\tableofcontents
\vspace{1em}

\noindent
This guide explains how to download and query each of the five materials databases
supported in the EuMINe DataBridge Hackathon.
You are required to use \textbf{at least two} of them in your submission.
All databases are publicly accessible free of charge; some require a (free) account
for API access.

\medskip
\noindent\textbf{Python environment required:}
\begin{lstlisting}[language=bash]
conda create -n eumine-hackathon python=3.11
conda activate eumine-hackathon
pip install mp-api pymatgen jarvis-tools requests pandas numpy
\end{lstlisting}

% ----------------------------------------------------------------------------
\section{Materials Project}

\subsection{Overview}
The Materials Project (\url{https://materialsproject.org}) provides DFT-computed
properties (PBE functional) for $\sim$150\,000 inorganic materials, including
formation energies, band gaps, elastic constants, and crystal structures.

\subsection{Getting an API key}
\begin{enumerate}
  \item Go to \url{https://materialsproject.org} and click \textbf{Sign in}.
  \item Create a free account (Google or email).
  \item Navigate to \textbf{Dashboard} $\to$ \textbf{API key}.
  \item Copy your key (a 32-character string).
  \item Store it as an environment variable --- \emph{never} hard-code it in scripts:
\begin{lstlisting}[language=bash]
export MP_API_KEY="your_key_here"
# or add to ~/.bashrc / ~/.zshrc
\end{lstlisting}
\end{enumerate}

\subsection{Installation}
\begin{lstlisting}[language=bash]
pip install mp-api pymatgen
\end{lstlisting}

\subsection{Basic query examples}

\paragraph{Query formation energy and band gap (no structures):}
\begin{lstlisting}[language=Python]
import os
from mp_api.client import MPRester

with MPRester(os.environ["MP_API_KEY"]) as mpr:
    docs = mpr.materials.summary.search(
        has_props=["thermo", "electronic_structure"],
        fields=["material_id", "formula_pretty",
                "formation_energy_per_atom", "band_gap",
                "nsites", "nelements", "symmetry"],
        chunk_size=1000,
        num_chunks=None,
    )
    for doc in docs:
        if doc.formation_energy_per_atom is None or doc.band_gap is None:
            continue
        print(doc.material_id, doc.formation_energy_per_atom, doc.band_gap)
\end{lstlisting}

\paragraph{Fetch a single structure by material ID:}
\begin{lstlisting}[language=Python]
from mp_api.client import MPRester
import os

with MPRester(os.environ["MP_API_KEY"]) as mpr:
    structure = mpr.get_structure_by_material_id("mp-22862")  # NaCl
print(structure)   # pymatgen Structure object
structure.to(fmt="cif", filename="NaCl.cif")
\end{lstlisting}

\paragraph{Batch fetch structures for a list of IDs:}
\begin{lstlisting}[language=Python]
from mp_api.client import MPRester
import os

ids = ["mp-22862", "mp-13", "mp-2657"]   # NaCl, Fe, TiO2

with MPRester(os.environ["MP_API_KEY"]) as mpr:
    structures = mpr.get_structures(ids)

for mid, struct in zip(ids, structures):
    struct.to(fmt="cif", filename=f"{mid}.cif")
\end{lstlisting}

\subsection{Rate limits and quotas}
\begin{itemize}
  \item Free tier: \textbf{30 requests/minute}, no daily cap.
  \item Use \verb|chunk_size=1000| and \verb|num_chunks=None| to stream large queries
        without hitting per-request limits.
  \item Avoid fetching \verb|structure| in bulk summary queries --- it multiplies
        bandwidth by 100$\times$. Fetch structures separately for only the materials
        you need.
\end{itemize}

% ----------------------------------------------------------------------------
\section{JARVIS-DFT}

\subsection{Overview}
JARVIS-DFT (\url{https://jarvis.nist.gov}) provides DFT-computed (OptB88vdW
functional) properties for $\sim$75\,000 materials with strong coverage of
2D and van-der-Waals materials.
Data is distributed as bulk JSON downloads via Figshare
--- no API key required.

\subsection{Installation}
\begin{lstlisting}[language=bash]
pip install jarvis-tools
\end{lstlisting}

\subsection{Bulk download (recommended)}
\begin{lstlisting}[language=Python]
from jarvis.db.figshare import data as jdata

# Downloads ~400 MB once; cached in ~/.jarvis/ for all subsequent calls
dft_3d = jdata("dft_3d")
print(f"Total entries: {len(dft_3d)}")   # ~75 000

for entry in dft_3d:
    jid      = entry.get("jid", "")
    formula  = entry.get("formula", "")
    ef       = entry.get("formation_energy_peratom")
    bg       = entry.get("optb88vdw_bandgap")
    natoms   = entry.get("natoms")
    spg      = entry.get("spg_number") or entry.get("spacegroup_number")
    print(jid, formula, ef, bg)
\end{lstlisting}

\subsection{Available datasets}
Other datasets accessible via the same interface:
\begin{center}
\begin{tabular}{ll}
\toprule
Dataset name & Description \\
\midrule
\texttt{dft\_3d}   & 3D bulk materials (main dataset) \\
\texttt{dft\_2d}   & 2D monolayers \\
\texttt{cfid\_3d}  & Classical force-field inspired descriptors \\
\texttt{jff}       & JARVIS force-field data \\
\bottomrule
\end{tabular}
\end{center}

\subsection{Accessing structures}
JARVIS entries include atomic positions as a nested dictionary compatible with
\verb|jarvis.core.atoms.Atoms|:
\begin{lstlisting}[language=Python]
from jarvis.core.atoms import Atoms
from jarvis.io.vasp.inputs import Poscar

entry = dft_3d[0]
atoms = Atoms.from_dict(entry["atoms"])
Poscar(atoms).write_file("structure.vasp")   # write as POSCAR
\end{lstlisting}

\paragraph{Convert to pymatgen Structure:}
\begin{lstlisting}[language=Python]
from pymatgen.io.vasp import Poscar as pmgPoscar

struct = pmgPoscar.from_file("structure.vasp").structure
struct.to(fmt="cif", filename="structure.cif")
\end{lstlisting}

% ----------------------------------------------------------------------------
\section{AFLOW}

\subsection{Overview}
AFLOW (\url{https://aflow.org}) is one of the largest repositories of
high-throughput DFT calculations, with $\sim$3.5 million entries across
multiple DFT flavors. Access is via a REST API --- no account required.

\subsection{REST API base URL}
\begin{center}
\url{http://aflow.org/API/aflowlib.json?}
\end{center}

\subsection{Query syntax}
Parameters are appended as \texttt{key=value} pairs separated by \texttt{\&}.

\paragraph{Retrieve entries with both formation energy and band gap:}
\begin{lstlisting}[language=bash]
curl "http://aflow.org/API/aflowlib.json?catalog=icsd&paging=1&\
Egap(*),enthalpy_formation_atom(*),compound(*),auid(*)&format=json" \
  | python -m json.tool | head -80
\end{lstlisting}

\paragraph{Python --- paginated download:}
\begin{lstlisting}[language=Python]
import requests, json, time

BASE = "http://aflow.org/API/aflowlib.json"
FIELDS = "auid,compound,enthalpy_formation_atom,Egap,spacegroup_relax"
PAGE_SIZE = 200

results = []
page = 1
while True:
    url = (f"{BASE}?catalog=icsd&paging={page}&"
           f"{FIELDS.replace(',', '(*),') + '(*)' }&format=json")
    r = requests.get(url, timeout=60)
    if r.status_code != 200 or not r.text.strip():
        break
    data = r.json()
    if not data:
        break
    results.extend(data)
    page += 1
    time.sleep(0.5)   # be polite --- 2 req/s is safe

print(f"Retrieved {len(results)} AFLOW entries")
\end{lstlisting}

\subsection{Key field names}
\begin{center}
\begin{tabular}{lll}
\toprule
AFLOW field & Property & Unit \\
\midrule
\texttt{enthalpy\_formation\_atom} & Formation energy per atom & eV/atom \\
\texttt{Egap}                      & Band gap (GGA)            & eV \\
\texttt{Egap\_type}                & Gap type (direct/indirect) & --- \\
\texttt{compound}                  & Chemical formula           & --- \\
\texttt{spacegroup\_relax}         & Space group number (relaxed) & --- \\
\texttt{natoms}                    & Number of atoms in cell    & --- \\
\bottomrule
\end{tabular}
\end{center}

\subsection{CIF download}
\begin{lstlisting}[language=Python]
import requests

auid = "aflow:some_auid_string"
# Replace : with / in AUID to build the file path
path = auid.replace("aflow:", "").replace(":", "/")
url  = f"http://aflow.org/{path}?files=CONTCAR.relax.vasp"
r = requests.get(url, timeout=30)
with open("structure.vasp", "wb") as f:
    f.write(r.content)
\end{lstlisting}

\subsection{Rate limits}
AFLOW has no documented rate limit but imposes soft throttling.
Keep requests below \textbf{2/second}; use \verb|time.sleep(0.5)| between calls.

% ----------------------------------------------------------------------------
\section{NOMAD}

\subsection{Overview}
NOMAD (\url{https://nomad-lab.eu}) aggregates DFT and experimental data from many
codes and groups worldwide, with rich metadata and a powerful REST API.
No account required for read access; a (free) NOMAD account enables uploads.

\subsection{REST API base URL}
\begin{center}
\url{https://nomad-lab.eu/prod/v1/api/v1/}
\end{center}

\paragraph{OpenAPI documentation:}
\url{https://nomad-lab.eu/prod/v1/api/v1/extensions/docs}

\subsection{Query for formation energy (Python)}
\begin{lstlisting}[language=Python]
import requests, json

BASE = "https://nomad-lab.eu/prod/v1/api/v1"

query = {
    "query": {
        "and": [
            {"results.material.structural_type": "bulk"},
            {"results.properties.electronic.band_structure_electronic:exists": True},
        ]
    },
    "required": {
        "results": {
            "material": {"chemical_formula_reduced": "*",
                         "symmetry": {"space_group_number": "*"}},
            "properties": {
                "thermodynamic": {"formation_energy": "*"},
                "electronic":    {"band_gap": "*"},
            }
        }
    },
    "pagination": {"page_size": 100}
}

r = requests.post(f"{BASE}/entries/query", json=query, timeout=60)
data = r.json()

for entry in data.get("data", []):
    mat  = entry["results"]["material"]
    prop = entry["results"]["properties"]
    ef   = prop.get("thermodynamic", {}).get("formation_energy")
    bg   = prop.get("electronic", {}).get("band_gap")
    formula = mat.get("chemical_formula_reduced")
    spg     = mat.get("symmetry", {}).get("space_group_number")
    print(formula, spg, ef, bg)
\end{lstlisting}

\subsection{Pagination}
NOMAD uses cursor-based pagination. To iterate all pages:
\begin{lstlisting}[language=Python]
all_entries = []
query["pagination"] = {"page_size": 100}
next_url = f"{BASE}/entries/query"

while next_url:
    r = requests.post(next_url, json=query, timeout=60)
    data = r.json()
    all_entries.extend(data.get("data", []))
    next_page = data.get("pagination", {}).get("next_page_after_value")
    if next_page:
        query["pagination"]["page_after_value"] = next_page
    else:
        break
\end{lstlisting}

\subsection{Key field names}
\begin{center}
\begin{tabular}{ll}
\toprule
NOMAD path & Property \\
\midrule
\texttt{results.properties.thermodynamic.formation\_energy} & Formation energy (eV/atom) \\
\texttt{results.properties.electronic.band\_gap}            & Band gap (eV) \\
\texttt{results.material.chemical\_formula\_reduced}        & Reduced formula \\
\texttt{results.material.symmetry.space\_group\_number}     & Space group \\
\bottomrule
\end{tabular}
\end{center}

\subsection{Rate limits}
NOMAD enforces \textbf{10 req/s} for anonymous users; use \verb|page_size=100| to
minimise the number of requests.

% ----------------------------------------------------------------------------
\section{OQMD}

\subsection{Overview}
The Open Quantum Materials Database (\url{https://oqmd.org}) contains $\sim$1 million
DFT entries computed with VASP/PBE.
Access is via a public REST API --- no account required.

\subsection{REST API base URL}
\begin{center}
\url{http://oqmd.org/oqmdapi/formationenergy?}
\end{center}

\subsection{Query examples}

\paragraph{Filter by element and request key properties:}
\begin{lstlisting}[language=bash]
curl "http://oqmd.org/oqmdapi/formationenergy?\
filter=element_set=(Ti,O)&\
fields=name,delta_e,band_gap,spacegroup,natoms&\
limit=50&offset=0&format=json"
\end{lstlisting}

\paragraph{Python --- paginated download:}
\begin{lstlisting}[language=Python]
import requests, time

BASE   = "http://oqmd.org/oqmdapi/formationenergy"
PARAMS = {
    "fields":  "name,delta_e,band_gap,spacegroup,natoms,unit_cell",
    "limit":   100,
    "offset":  0,
    "format":  "json",
}

results = []
while True:
    r = requests.get(BASE, params=PARAMS, timeout=60)
    data = r.json()
    batch = data.get("data", [])
    if not batch:
        break
    results.extend(batch)
    PARAMS["offset"] += PARAMS["limit"]
    time.sleep(0.5)
    if PARAMS["offset"] > data.get("meta", {}).get("count", 0):
        break

print(f"Retrieved {len(results)} OQMD entries")
\end{lstlisting}

\subsection{Key field names}
\begin{center}
\begin{tabular}{lll}
\toprule
OQMD field     & Property             & Unit \\
\midrule
\texttt{name}      & Chemical formula         & --- \\
\texttt{delta\_e}  & Formation enthalpy/atom  & eV/atom \\
\texttt{band\_gap} & Band gap                 & eV \\
\texttt{spacegroup}& Space group symbol       & --- \\
\texttt{natoms}    & Atoms in cell            & --- \\
\bottomrule
\end{tabular}
\end{center}

\subsection{Rate limits}
OQMD has no documented rate limit. Keep requests below \textbf{1/second} to avoid
triggering the server's soft throttle.

% ----------------------------------------------------------------------------
\section{Common Pitfalls and Best Practices}

\subsection{API rate limits}
\begin{itemize}
  \item \textbf{Materials Project:} 30 req/min --- use the official client
        (\texttt{mp-api}) which handles chunking automatically.
  \item \textbf{JARVIS:} no API; single bulk download. Run it once, cache the result.
  \item \textbf{AFLOW:} no documented limit; stay at $\leq$2 req/s.
  \item \textbf{NOMAD:} 10 req/s; always use large \texttt{page\_size}.
  \item \textbf{OQMD:} polite limit; 1 req/s recommended.
\end{itemize}

\subsection{Pagination}
All REST APIs return a limited number of results per call.
Always implement pagination and check for empty responses.
A common pattern:
\begin{lstlisting}[language=Python]
offset = 0
while True:
    batch = fetch_page(offset=offset, limit=PAGE_SIZE)
    if not batch:
        break
    process(batch)
    offset += PAGE_SIZE
\end{lstlisting}

\subsection{Property name differences}
The same quantity has different field names across databases.
Use the following mapping when harmonising data:

\begin{center}
\begin{tabular}{llll}
\toprule
Database & Formation energy field & Band gap field & Unit \\
\midrule
Materials Project & \texttt{formation\_energy\_per\_atom} & \texttt{band\_gap} & eV, eV \\
JARVIS-DFT        & \texttt{formation\_energy\_peratom}  & \texttt{optb88vdw\_bandgap} & eV, eV \\
AFLOW             & \texttt{enthalpy\_formation\_atom}   & \texttt{Egap} & eV, eV \\
NOMAD             & \texttt{formation\_energy}           & \texttt{band\_gap} & eV/atom, eV \\
OQMD              & \texttt{delta\_e}                    & \texttt{band\_gap} & eV/atom, eV \\
\bottomrule
\end{tabular}
\end{center}

\noindent\textbf{Note:} Always check units per database.
NOMAD's \texttt{formation\_energy} is sometimes in eV for the whole cell, not per atom
--- verify by dividing by \texttt{natoms} and cross-checking with a known compound.

\subsection{Missing and invalid values}
Databases use different sentinels for missing data:
\begin{lstlisting}[language=Python]
MISSING = {None, "na", "Na", "NA", "", "None", "N/A", float("nan")}

def is_valid(value):
    if value in MISSING:
        return False
    try:
        return not (float(value) != float(value))   # NaN check
    except (TypeError, ValueError):
        return False
\end{lstlisting}

\subsection{Functional and methodological biases}
The same compound can yield very different values in different databases because the
DFT functional, pseudopotentials, and convergence settings differ:

\begin{center}
\begin{tabular}{ll}
\toprule
Database & Primary functional \\
\midrule
Materials Project & PBE (GGA) \\
JARVIS-DFT        & OptB88vdW (vdW-corrected GGA) \\
AFLOW             & PBE (GGA), LDA for some \\
NOMAD             & Mixed (source-dependent) \\
OQMD              & PBE (GGA) \\
\bottomrule
\end{tabular}
\end{center}

Band gaps are \textbf{systematically underestimated} by standard GGA.
JARVIS reports slightly higher band gaps than MP for the same compound due to the
OptB88vdW correction, especially for weakly bonded and layered systems.
Document which sources you used and how you resolved discrepancies in your
\textit{Data Integration Report}.

\subsection{Formula and structure matching}
When linking entries across databases:
\begin{itemize}
  \item Normalize all formulas to \textbf{reduced composition}
        (e.g., \texttt{Na2Cl2} $\to$ \texttt{NaCl}).
        Use \texttt{pymatgen.core.Composition(formula).reduced\_formula}.
  \item Match on \textbf{reduced formula + space group number} as a primary key.
  \item Apply an \textbf{nsites tolerance of $\pm2$} to reconcile slight
        supercell differences.
  \item Deduplicate after merging (keep one MP entry and one JARVIS entry per
        bridge material to avoid data leakage).
\end{itemize}

\subsection{Storing large datasets efficiently}
\begin{itemize}
  \item Use \textbf{JSONL} (one JSON object per line) for streaming writes ---
        safe against RAM overflow and easy to append to.
  \item Avoid loading all 150,000 MP entries into a Python list at once.
  \item For the hackathon, the provided Bridge Dataset ($\sim$1,000 entries) is
        small enough to hold in memory as a Pandas DataFrame.
\end{itemize}

% ----------------------------------------------------------------------------
\section*{Summary: Quick Reference}
\addcontentsline{toc}{section}{Summary: Quick Reference}

\begin{center}
\begin{tabular}{lllll}
\toprule
Database & Size & API type & Auth & Install \\
\midrule
Materials Project & $\sim$150k & Python client & API key (free) & \texttt{pip install mp-api} \\
JARVIS-DFT        & $\sim$75k  & Bulk download & None           & \texttt{pip install jarvis-tools} \\
AFLOW             & $\sim$3.5M & REST/JSON      & None           & \texttt{pip install requests} \\
NOMAD             & Diverse    & REST/JSON      & None (read)    & \texttt{pip install requests} \\
OQMD              & $\sim$1M   & REST/JSON      & None           & \texttt{pip install requests} \\
\bottomrule
\end{tabular}
\end{center}

\vspace{2em}
\noindent\rule{\linewidth}{0.4pt}
{\small
\textit{EuMINe DataBridge Hackathon 2026} \hfill \textit{hackathon@eumine.eu}\\
Questions about data access? Post in \texttt{\#data-sources} on Slack or open a
GitHub Discussion at
\url{https://github.com/EuMINe-COST/eumine_hackathon_2026/discussions}.
}

\end{document}
